%% LyX 2.4.0~RC4 created this file.  For more info, see https://www.lyx.org/.
%% Do not edit unless you really know what you are doing.
\documentclass[letterpaper]{article}
\usepackage{lmodern}
\renewcommand{\sfdefault}{lmss}
\renewcommand{\ttdefault}{lmtt}
\usepackage[T1]{fontenc}
\usepackage[latin9]{inputenc}
\synctex=-1
\usepackage{refstyle}
\usepackage{float}
\usepackage{enumitem}
\usepackage{amsmath}
\usepackage{amssymb}
\usepackage[pdfusetitle,
 bookmarks=true,bookmarksnumbered=false,bookmarksopen=false,
 breaklinks=false,pdfborder={0 0 1},backref=false,colorlinks=false]
 {hyperref}

\makeatletter

%%%%%%%%%%%%%%%%%%%%%%%%%%%%%% LyX specific LaTeX commands.

\global\long\def\N{\mathbb{N}}%
\global\long\def\Z{\mathbb{Z}}%
\global\long\def\Q{\mathbb{Q}}%


\global\long\def\Lang{\mathcal{L}}%


\global\long\def\iswd{\mathord{\downarrow}}%


\global\long\def\isnotwd{\mathord{\uparrow}}%


\global\long\def\turing{\mathrm{T}}%


\global\long\def\upto{\mathord{\upharpoonright}}%
\global\long\def\ZFniptc{\mathsf{ZF}\mathord{-}\mathsf{inf}\mathord{+}\mathsf{TC}}%
\global\long\def\ZFfinptc{\mathsf{ZFfin}\mathord{+}\mathsf{TC}}%
\global\long\def\restrict#1{\left.#1\right|}%
\global\long\def\PP{\mathcal{P}}%
\global\long\def\PPfin{\PP^{\mathrm{fin}}}%


\global\long\def\llor{\bigvee\mathchoice{\hspace{-1.1em}}{\hspace{-0.7em}}{\hspace{-0.7em}}{\hspace{-0.6em}}\bigvee}%
\global\long\def\lland{\bigwedge\mathchoice{\hspace{-1.1em}}{\hspace{-0.7em}}{\hspace{-0.7em}}{\hspace{-0.6em}}\bigwedge}%
\global\long\def\llorc{\llor^{\mathrm{c}}}%
\global\long\def\eqvrel{\simeq}%
\global\long\def\etau{\varepsilon\tau}%


\AtBeginDocument{\providecommand\defref[1]{\ref{def:#1}}}
\AtBeginDocument{\providecommand\thmref[1]{\ref{thm:#1}}}
\AtBeginDocument{\providecommand\secref[1]{\ref{sec:#1}}}
\AtBeginDocument{\providecommand\rmkref[1]{\ref{rmk:#1}}}
\AtBeginDocument{\providecommand\enuref[1]{\ref{enu:#1}}}
\AtBeginDocument{\providecommand\lemref[1]{\ref{lem:#1}}}
\AtBeginDocument{\providecommand\propref[1]{\ref{prop:#1}}}
\pdfpageheight\paperheight
\pdfpagewidth\paperwidth

\RS@ifundefined{subsecref}
  {\newref{subsec}{name = \RSsectxt}}
  {}
\RS@ifundefined{thmref}
  {\def\RSthmtxt{theorem~}\newref{thm}{name = \RSthmtxt}}
  {}
\RS@ifundefined{lemref}
  {\def\RSlemtxt{lemma~}\newref{lem}{name = \RSlemtxt}}
  {}


%%%%%%%%%%%%%%%%%%%%%%%%%%%%%% Textclass specific LaTeX commands.
\usepackage[amsmath,thmmarks]{ntheorem}
\newlength{\lyxlabelwidth}      % auxiliary length 
\newref{thm}{refcmd={Theorem~\ref{#1}}}
\theoremstyle{plain}
\theorembodyfont{\upshape}
\theoremheaderfont{\bfseries}
\theoremseparator{.}
\theoremnumbering{arabic}
\theoremsymbol{}
\newtheorem{thm}{Theorem}[subsection]
\newref{def}{refcmd={Definition~\ref{#1}}}
		\newtheorem{definition}[thm]{Definition}
\newref{rmk}{refcmd={Remark~\ref{#1}}}
		\newtheorem{remark}[thm]{Remark}
\theoremstyle{nonumberplain}
\theoremheaderfont{\itshape}
\theorembodyfont{\upshape}
\theoremseparator{:}
\theoremsymbol{\ensuremath{\blacksquare}}
\newtheorem{proof}{Proof}

\theoremstyle{plain}
\theorembodyfont{\upshape}
\theoremheaderfont{\bfseries}
\theoremseparator{.}
\theoremnumbering{arabic}
\theoremsymbol{}
\theoremstyle{nonumberplain}
\theorembodyfont{\upshape}
\theoremheaderfont{\bfseries}
\theoremseparator{.}
\newtheorem{question**}{Question}
\theoremstyle{plain}
\theorembodyfont{\upshape}
\theoremheaderfont{\bfseries}
\theoremseparator{.}
\theoremnumbering{arabic}
\theoremsymbol{}
\theoremstyle{nonumberplain}
\theorembodyfont{\upshape}
\theoremheaderfont{\bfseries}
\theoremseparator{.}
\renewtheorem{thm*}{Theorem}
\theoremstyle{plain}
\theorembodyfont{\upshape}
\theoremheaderfont{\bfseries}
\theoremseparator{.}
\theoremnumbering{arabic}
\theoremsymbol{}
\newref{prop}{refcmd={Proposition~\ref{#1}}}
		\newtheorem{prop}[thm]{Proposition}
\newref{lem}{refcmd={Lemma~\ref{#1}}}
		\newtheorem{lemma}[thm]{Lemma}
\newref{question}{refcmd={Question~\ref{#1}}}
		\newtheorem{question}{Question}

%%%%%%%%%%%%%%%%%%%%%%%%%%%%%% User specified LaTeX commands.

\usepackage{listings}
\lstset{
basicstyle=\ttfamily\footnotesize,
keepspaces=true,
tabsize=4,
breaklines=true,
columns=fullflexible,
mathescape=true
}
\usepackage{fullpage}
\usepackage{braket}


\usepackage[cal=euler]{mathalpha}

\setlist[enumerate,2]{ref=\theenumi.\theenumii}
\setlist[enumerate,3]{ref=\theenumi.\theenumii.\theenumiii}
\setlist[enumerate,4]{ref=\theenumi.\theenumii.\theenumiii.\theenumiv}


\newref{sec}{refcmd={Section \ref{#1}}}

\usepackage[style=numeric]{biblatex}
\DeclareFieldFormat{postnote}{#1}

\usepackage{orcidlink}

\makeatother

\usepackage[style=numeric]{biblatex}
\addbibresource{bibliography.bib}
\begin{document}
\title{Escaping Tennenbaum's Theorem and a Strong Jump Inversion Theorem}
\author{Duarte Maia \orcidlink{0009-0004-8658-4840}}

\maketitle

\section{Introduction}

\subsection{Background -- Tennenbaum's Theorem}

\nocite{lutz_walsh_tennenbaum}

It is a well-known result of Tennenbaum (see e.g.\ \cite{kaye_book})
that $\mathsf{PA}$, the first-order theory of Peano Arithmetic, does
not admit any non-standard computable model. In fact, in every non-standard
model of $\mathsf{PA}$, neither addition nor multiplication is computable.
Generalizations of this theorem can be found in the literature, including
versions of it for weaker fragments of $\mathsf{PA}$ \cite{wilmers_PA},
alternative operations \cite{schmerl_tennenbaum_recursive_reducts},
and finite set theory \cite{mancini_zambella_set_theories,computable_quotient_arithmetic_set_theory},
to name a few. However, all of these approaches and variations appear
to depend heavily on the specificities of the chosen signature. This
was observed by Pakhomov, who in 2022 \cite{pakhomov} constructed
a theory that is definitionally equivalent to $\mathsf{PA}$ (see
\defref{defeqv} below) for which there is a computable nonstandard
model. This shows that Tennenbaum's theorem really is reliant on the
choice of signature.

In this paper, we extend Pakhomov's construction to the stronger theories
obtained by adding to $\mathsf{PA}$ all true $\Pi_{n}$ sentences,
for fixed values of $n$ (\thmref{pakhomov_answer}). In the process,
we obtain a general theorem for strong jump inversion (\thmref{main_final}),
and show that this theorem can be used to establish multiple other
results already present in the literature (\secref{prevs_applications}).
\begin{definition}[Definitional Extension]
Let $T$ be a theory over a signature $\Lang$. A \emph{definitional
extension} of $T$ is a theory $S\supseteq T$, over a signature $\Lang'\supseteq\Lang$,
such that
\begin{itemize}
\item Every theorem of $S$ that only uses symbols from the signature $\Lang$
is also a theorem of $T$, and
\item Every symbol of $\Lang'\setminus\Lang$ is $S$-definable in terms
of symbols in $\Lang$. For example, if $P$ is any predicate symbol
in $\Lang'$, there is an $\Lang$-formula $\varphi$ such that
\[
S\vdash(\forall\vec{x})(P(\vec{x})\leftrightarrow\varphi(\vec{x})),
\]
and a similar statement holds for function symbols.
\end{itemize}
\end{definition}
%
\begin{definition}[Definitional Equivalence]
\label{def:defeqv}Let $T_{1}$ and $T_{2}$ be two theories, whose
signatures $\Lang_{1}$ and $\Lang_{2}$ respectively are assumed
to be disjoint with no loss of generality. We say that $T_{1}$ and
$T_{2}$ are \emph{definitionally equivalent} if there is a theory
$T$ which is a definitional extension of both.
\end{definition}
\begin{remark}
The notion of definitional equivalence is related, but not identical,
to a more common notion of ``equal power of theories'' called bi-interpretability:
definitional equivalence is strictly stronger, though these two notions
coincide for most natural examples. The interested reader will find
a thorough comparison of these two notions in \cite{biinterpretabilityvssynonymy}.
\end{remark}
A nontrivial example of two definitionally equivalent theories is
that of $\mathsf{PA}$ and a specific version of finite set theory,
which we will refer to as $\ZFfinptc$. This theory, and a closely
related one, will play an important role in what follows, so we formally
introduce them to the reader. A study of this theory can be found
in \cite{kaye_wong}.
\begin{definition}
The theory $\ZFniptc$ consists of the usual axioms of Zermelo-Fraenkel
set theory, with the removal of the axiom of infinity, and with the
addition of the so-called ``axiom of transitive closure'', which
for our purposes is equivalent to the statement: $V=\cup_{\alpha\in\mathrm{Ord}}V_{\alpha}$.
\end{definition}
%
\begin{definition}
The theory $\ZFfinptc$ consists of $\ZFniptc$, together with the
negation of the axiom of infinity.
\end{definition}
\begin{thm}
\label{thm:pazffineqv}The theories $\mathsf{PA}$ and $\ZFfinptc$
are definitionally equivalent.
\end{thm}
\begin{proof}
See \cite{kaye_wong}.
\end{proof}

\subsection{Pakhomov's Construction and New Results}

In \cite{pakhomov}, Pakhomov constructed a theory that is definitionally
equivalent to $\ZFniptc$, and consequently to $\mathsf{PA}$, and
has computable nonstandard models. We start by reviewing his results.
\begin{definition}
\label{def:tboth}Let $T(S)$ be the theory defined by Pakhomov, which
axiomatizes a single ternary predicate $S(x,y,z)$. The theory $T(S)$
is definitionally equivalent to $\ZFniptc$. We let $T(\in,S)$ be
the common definitional extension.
\end{definition}
\begin{thm}[Pakhomov, \cite{pakhomov}]
\label{thm:premain}If $T$ is a consistent c.e.\ theory extending
$T(S)$, then $T$ admits a computable model.
\end{thm}
Pakhomov proves this using an explicit ``Henkin construction'' type
of argument. We were able to reframe his work in terms of the following
stronger theorem:
\begin{thm}
If $D$ is a $0'$-computable model of $T(\in,S)$, the reduct $D\upto S$
admits a computable copy.
\end{thm}
\begin{proof}
Case $n=1$ of \thmref{nestedpakhomov} below.
\end{proof}
In this paper, we generalize the methods used by Pakhomov to answer
a question posed by him in the affirmative:
\begin{question**}[Pakhomov \cite{pakhomov}]
Fix a value of $n$. Are there theories definitionally equivalent
to ``$\mathsf{PA}$ plus all true $\Pi_{n}$ sentences'' that have
computable non-standard models?
\end{question**}
The motivation for this question is that, as Pakhomov proved in \cite{pakhomov},
this statement is not true `in the limit'. More precisely, Pakhomov
showed that any nonstandard model of a theory that is definitionally
equivalent to true arithmetic cannot be computable. This led to the
question of whether a partial result could be recovered.

We answer Pakhomov's question via the following improvement on Pakhomov's
construction of $T(S)$:
\begin{thm*}[\ref{thm:nestedpakhomov}]
There is a nested sequence of consistent c.e.\ theories $\ZFniptc\subseteq T^{0}\subseteq T^{1}\subseteq T^{2}\subseteq\cdots$
satisfying the following properties:
\begin{itemize}
\item Each $T^{n}$ is in the signature containing $\in$ and predicates
$S^{0}$, $\dots$, $S^{n}$, with each $S^{i}$ being an $(i+2)$-ary
predicate symbol,
\item All of these extensions are conservative over $\ZFniptc$, in the
sense that they contain no additional theorems in the predicate $\in$,
\item Each of the predicates $\in$ and $S^{n}$ is definable from the other
within $T^{n}$, and
\item Given an $X'$-computable model $D$ of $T^{n}$ restricted to the
signature containing $S^{i}$, $\dots$, and $S^{n}$, there is an
$X$-computable copy $M$ of $D\upto(S^{i+1},\dots,S^{n})$.
\end{itemize}
\end{thm*}
\begin{remark}
The second and third bullet point imply that, for every $n$, $\ZFniptc$
and $T^{n}\upto S^{n}$ are definitionally equivalent.
\end{remark}
From \thmref{nestedpakhomov}, we conclude one of the main results
of this paper:
\begin{thm*}[\ref{thm:pakhomov_answer}]
For every $n$, there is a theory definitionally equivalent to ``$\mathsf{PA}$
plus all true $\Pi_{n}$ sentences'' that admits a computable non-standard
model.
\end{thm*}
\begin{proof}
First, note that the theory $T$ containing $\mathsf{PA}$ plus all
true $\Pi_{n}$ sentences is a $0^{(n)}$-c.e.\ theory. By the results
of Kaye and Wong \cite{kaye_wong}, $\mathsf{PA}$ is definitionally
equivalent to $\ZFfinptc$, and therefore $T$ is definitionally equivalent
to a $0^{(n)}$-c.e.\ extension of the theory $T^{n+1}$ from \thmref{nestedpakhomov}.
Say $\bar{T}$ is this theory, and let $\bar{T}^{+}$ be the same
with an added constant $c$ and an axiom schema ensuring that $c$
is a nonstandard element. Let $D$ be a $0^{(n+1)}$-computable model
of $\bar{T}^{+}$, in the signature containing the $n+2$ predicates
$S^{0}$ to $S^{n+1}$, which exists by the computable completeness
theorem (see \cite{harizanov}). Then, apply the relativized version
of \thmref{nestedpakhomov} $n+1$ times, obtaining a computable model
of $\bar{T}\upto S^{n+1}$. This shows that $\bar{T}\upto S^{n+1}$
is the theory we sought: It is definitionally equivalent to ``$\mathsf{PA}$
plus all true $\Pi_{n}$ sentences'', and it admits a computable
nonstandard model.
\end{proof}
Pakhomov \cite{pakhomov} also posed another question, which has been
answered by Lutz and Walsh.
\begin{question**}[Pakhomov \cite{pakhomov}]
Is there a c.e.\ theory $T$ such that no definitionally equivalent
theory $T'$ admits a computable model?
\end{question**}
\begin{thm}[Lutz--Walsh \cite{lutz_walsh_tennenbaum}]
There is a c.e.\ theory $T$ such that every model of every theory
$T'$ definitionally equivalent to $T$ is noncomputable.
\end{thm}

\subsection{A Strong Jump Inversion Theorem}

In the process of answering Pakhomov's question, it happened that
certain parts of our construction were not specific to arithmetic,
and so we were able to extract a general-purpose theorem for strong
jump inversion.\footnote{Patrick Lutz was the first person to bring the idea of a broad generalization
to my attention. He also provided a working conjecture that was very
useful in developing the theory.}
\begin{remark}
For general background on Strong Jump Inversion theorems and their
history, see Section 1.1 of \cite{strong_jump_inversion}. A general
reference for computable structure theory is \cite{montalban_cst}.
\end{remark}
%
\begin{remark}
We are not the first to seek a general-purpose jump inversion theorem.
The 2018 paper \cite{strong_jump_inversion} also contains a general-purpose
theorem (Theorem 2.5) from which some theorems from the literature
are recovered. Our theorem seems to follow similar themes as the theorem
found in \cite{strong_jump_inversion}, but our assumptions have a
different flavor and our conclusion goes in a different direction.
Notably, our result requires some amount of creativity in its application
-- one must, in most cases, construct an intermediate structure to
apply our \thmref{main_final} to  obtain a jump inversion result,
while Theorem 2.5 from \cite{strong_jump_inversion} is relatively
straight-forward in its application to examples.
\end{remark}
To present our main result, \thmref{main_final}, we need some definitions.
From this moment on, all structures are assumed to be relational.

\begin{definition}
\label{def:atomicformula}For us, an \emph{atomic formula} in the
variables $\vec{x}$ is a formula obtained by applying a single \emph{positive
}instance of a predicate to some of the variables in the tuple $\vec{x}$
(possibly repeated and reordered). We do not consider the negation
of a predicate to be atomic.
\end{definition}
%
\begin{definition}
\label{def:atomictype}An \emph{atomic type} in the variables $\vec{x}$
is a set of atomic formulas in $\vec{x}$. For a tuple of elements
$\vec{b}$ in a structure, the \emph{atomic type of $\vec{b}$} consists
of the set of atomic formulas satisfied by $\vec{b}$.
\end{definition}
%
\begin{definition}[$0'$-computable c.e.-typed structure]
\label{def:cetypedpresentation}Let $\Lang'$ be a computable relational
signature.
\global\long\def\tp{\mathrm{tp}}%
 A\emph{ $0'$-computable c.e.-typed structure} $D$ consists of a
$0'$-computable function $t$, whose domain is the finite power-set
of an initial segment of $\N$ (this initial segment we call the \emph{domain}
of $D$), which takes as input the strong index of a finite set of
natural numbers $\vec{n}=\{n_{0}<n_{1}<\dots<n_{k-1}\}$ and outputs
the c.e.\ index $t(\vec{n})$ for an atomic $\Lang'$-type $\tp(\vec{n})$
in the variables $(x_{n_{0}},\dots,x_{n_{k-1}})$, with the compatibility
condition that if $\vec{n}\subseteq\vec{m}$ then $\tp(\vec{n})\subseteq\tp(\vec{m})$.
\end{definition}
%
\begin{definition}
\label{def:comp_inf_lang}We work in the Computable Infinitary Language,
as described in \cites[Chapter III]{montalban_cst2}.

Given a relational signature $\Lang'$, we define a $\llorc_{0}$
formula as a formula of the type $\varphi(\vec{x})\equiv P_{1}\land\dots\land P_{n}$,
where each $P_{j}$ is an atomic formula in $\vec{x}$, per \defref{atomicformula}.
We also allow for $P_{j}$ to denote $\top$, the ``true'' predicate,
and $\bot$, the ``false'' predicate.

We define a $\llorc_{1}$ formula as a formula (in finitely many variables)
of the type
\[
\varphi(\vec{x})\equiv\llor_{i<\omega}\varphi_{i}(\vec{x}),
\]
where $\{\varphi_{i}\}_{i<\omega}$ is a computable sequence of $\llorc_{0}$
formulas.

Finally, we define the $\llorc_{2}$ formulas as those formulas in
finitely many variables of the type
\[
\varphi(\vec{x})\equiv\llor_{i<\omega}\left(\psi_{i}(\vec{x})\land\neg\varphi_{i}(\vec{x})\right),
\]
where $\{\psi_{i}\}_{i<\omega}$ and $\{\varphi_{i}\}_{i<\omega}$
are computable sequences of $\llorc_{1}$ formulas.

One might imagine extending the definition to higher levels, by defining
a $\llorc_{n+1}$ formula as an infinite c.e.\ disjunction of conjunctions
of $\llorc_{n}$ formulas and negations thereof, but these first levels
will be sufficient for our purposes.
\end{definition}
%
\begin{definition}[QETP]
\label{def:qetp}Let $\Lang\subseteq\Lang'$ be computable signatures,
and let $D$ be a structure over $\Lang'$. We say that $D$ satisfies
the \emph{Computable Positive Quantifier Elimination and Trash Existence
Property}, abbreviated to \emph{QETP,} if the following properties
hold:
\begin{enumerate}[label=(\alph*)]
\item \label{enu:qetp0}$\Lang$ is finite, closed under negation, and
includes the predicate $\neq$,
\item \label{enu:qetp1}For every quantifier-free $\Lang$-type $q(\vec{x},\vec{y})$,
there is a $\llorc_{1}$\emph{ }$\Lang'$-formula $\chi_{q}(\vec{x})$,
computable from $q$, such that
\[
D\vDash(\forall\vec{x})\left[\chi_{q}(\vec{x})\leftrightarrow(\exists\vec{y})\,q(\vec{x},\vec{y})\right],
\]
\item \label{enu:qetp2}There is a $0'$-computable \emph{partial} function,
denoted by $\tau$, which takes as input a pair $(q(\vec{x},y,\vec{z}),\varphi(\vec{x}))$,
where $q(\vec{x},y,\vec{z})$ is a quantifier-free $\Lang$-type\footnote{Note the lack of ``vector arrow'' on the variable $y$. This denotes
a \emph{single} variable, not a tuple of them.} and $\varphi(\vec{x})$ is a \emph{finitary} $\Lang'$-formula, whose
output is a $\llorc_{2}$ $\Lang'$-formula denoted $\tau_{q\varphi}(\vec{x},y)$
that satisfies the following two properties:
\begin{itemize}
\item For every tuple $(\vec{x},y,\vec{z})$ of elements in $D$, of $\Lang$-type
$q(\vec{x},y,\vec{z})$, there is a formula $\varphi(\vec{x})$ satisfied
by $\vec{x}$ such that $(q,\varphi)$ is in the domain of $\tau$.
Equivalently, if $\vec{b}$ is a tuple of elements in $D$ such that
$D\vDash\chi_{q}(\vec{b})$, there is $\varphi(\vec{x})$ such that
$(q,\varphi)$ is in the domain of $\tau$ and $D\vDash\varphi(\vec{b})$.
Moreover,
\item The following formula holds in $D$:
\[
D\vDash(\forall\vec{x})\left[\chi_{q}(\vec{x})\land\varphi(\vec{x})\rightarrow(\exists y)\tau_{q\varphi}(\vec{x},y)\right].
\]
\end{itemize}
\item \label{enu:qetpetau}Given a pair $(q(\vec{x},y,\vec{z}),\varphi(\vec{x}))$
in the domain of $\tau$ as above, together with a quantifier-free
$\Lang$-type $Q(\vec{x},y,\vec{z},\vec{w})$ extending $q(\vec{x},y,\vec{z})$,
there is a $\llorc_{1}$ $\Lang'$-formula $\etau_{q\varphi Q}(\vec{x})$,
computable from the triple $(q,\varphi,Q)$, such that
\[
D\vDash(\forall\vec{x})\biggl[(\exists y)\left[\chi_{Q}(\vec{x},y)\land\tau_{q\varphi}(\vec{x},y)\right]\leftrightarrow\etau_{q\varphi Q}(\vec{x})\biggr].
\]
The function $(q,\varphi,Q)\mapsto\etau_{q\varphi Q}$ may be defined
in some cases even if the pair $(q,\varphi)$ is not in the domain
of $\tau$. In this scenario, we do not make any requirement of $\etau_{q\varphi Q}$.

Note: As a corollary of \ref{enu:qetp2}, we obtain $D\vDash(\forall\vec{x})\left[\chi_{q}(\vec{x})\land\varphi(\vec{x})\rightarrow\etau_{q\varphi q}(\vec{x})\right]$
for $(q,\varphi)$ in the domain of $\tau$.
\item \label{enu:qetp3}Finally, for all $q$, $\varphi$, $Q$, we require
\[
D\vDash(\forall\vec{x},y)\left[\etau_{q\varphi Q}(\vec{x})\land\tau_{q\varphi}(\vec{x},y)\rightarrow\chi_{Q}(\vec{x},y)\right].
\]
\end{enumerate}
\end{definition}
We now state the central result of this paper.
\begin{thm*}[\ref{thm:main_final}]
If $D$ is a $0'$-computable c.e.-typed structure over the signature
$\Lang'\supseteq\Lang$ satisfying the QETP, then the reduct $D\upto\Lang$
admits a computable copy $M$. There is a $0'$-computable isomorphism
between $D\upto\Lang$ and $M$. This result is uniform in the index
for $D$ and an index for witnesses that $D$ satisfies the QETP,\footnote{This includes: An index for which predicates in $\Lang$ are the negation
of each other, an index for the computable functions $\chi$ and $\etau$,
and an index for the $0'$-computable function $\tau$.} and relativizes uniformly.
\end{thm*}
We obtain some known results as corollaries of \thmref{main_final}
-- see \secref{prevs_applications}.

An extended version of this article can be found in \cite{self_escaping_tennenbaum}.

\section{A General Jump Inversion Theorem}

\subsection{\protect\label{subsec:qetp}Preliminary Observations}

Before the main proof of this paper, we make some observations that
will be necessary regarding the new types of formulas we introduced
in \defref{comp_inf_lang} and regarding the definition of the QETP,
\defref{qetp}.

\begin{remark}
\label{rmk:llorc}The main property of the $\llorc_{1}$ formulas
is: Given a c.e.\ index for the atomic $\Lang'$-type of a tuple
of elements $\vec{b}$ in some structure, and a $\llorc_{1}$ formula
$\varphi(\vec{x})$ satisfied by $\vec{b}$, we can computably confirm
in finite time that the tuple $\vec{b}$ satisfies the formula $\varphi$.
\end{remark}
%
\begin{remark}
\label{rmk:llorc2}The main property of the $\llorc_{2}$ formulas
is: Given a tuple of elements $\vec{b}$ of a structure and a c.e.\ index
for the atomic type of $\vec{b}$, it is $0'$-c.e.\ to determine
whether a $\llorc_{2}$ formula holds of $\vec{b}$.

To elaborate: If $\varphi(\vec{x})\equiv\llor_{i}\left(\psi_{i}(\vec{x})\land\neg\varphi_{i}(\vec{x})\right)$
is a $\llorc_{2}$ formula, and $\vec{b}$ is a tuple of elements
of a structure for which we have a c.e.\ index for the atomic type
of $\vec{b}$, we can computably create a sequence of guesses for
whether $\varphi(\vec{b})$ holds: Guess `no' until we find $i$
for which $\psi_{i}(\vec{b})$ holds. Then guess `yes' until we
find that $\varphi_{i}(\vec{b})$ holds. Again guess `no' until
we find a new $j$ for which $\psi_{j}(\vec{b})$ holds, then `yes'
until we find that $\varphi_{j}(\vec{b})$ holds, and so on. This
sequence of guesses will have the following behavior:
\begin{itemize}
\item If $\varphi(\vec{b})$ holds, we will guess `yes' all but finitely
many times, and
\item If $\varphi(\vec{b})$ does not hold, we will guess `no' infinitely
many times.
\end{itemize}
\end{remark}
%
\begin{remark}
\label{rmk:fokinaetal}Bazhenov, Fokina, Rossegger, A.~Soskova, and
Vatev introduced in \cite{lopez_escobar_cts_domains} two hierarchies
of infinitary formulas, which they named $\Sigma_{\alpha}^{p}$ and
$\Pi_{\alpha}^{p}$ for $\alpha<\omega_{1}$, which are defined with
positivity of predicates in mind.\footnote{We would like to thank an anonymous reviewer for bringing this work
to our attention.} These are closely related to our work: Our $\llorc_{n}$ formulas
coincide exactly with the set of quantifier-free $\Sigma_{n}^{p}$
formulas.
\end{remark}
%
\begin{remark}
\label{rmk:tauwork}We will work with Item \ref{enu:qetp2} of the
QETP in the following manner: We envision a computable process that
enumerates triples $(q,\varphi,\tau)$, and sometimes deletes them,
in such a manner that, if $\tau_{q\varphi}$ is well-defined, the
triple $(q,\varphi,\tau_{q\varphi})$ will eventually be enumerated
and never removed, and every triple that is not of the form $(q,\varphi,\tau_{q\varphi})$
for $(q,\varphi)$ in the domain of $\tau$ will eventually be removed.
\end{remark}
%
\begin{remark}
\label{rmk:weakerchi}We make an additional technical assumption on
the formulas $\chi_{q}$ with no loss of generality. Let $q(\vec{x},\vec{y},\vec{z})$
be a quantifier-free $\Lang$ formula. We can assign a $\chi$-formula
to $q$ in (at least) two different ways: By grouping the variables
as $(\vec{x},[\vec{y},\vec{z}])$, or by grouping them as $([\vec{x},\vec{y}],\vec{z})$.
This leads to two different formulas, which, when necessary to distinguish,
we notate respectively as $\chi_{q(\vec{x},\cdot,\cdot)}(\vec{x})$
and $\chi_{q(\vec{x},\vec{y},\cdot)}(\vec{x},\vec{y})$. Now, \textbf{we
assume without loss of generality} that, in such a scenario, $\chi_{q(\vec{x},\vec{y},\cdot)}(\vec{x},\vec{y})$
contains $\chi_{q(\vec{x},\cdot,\cdot)}(\vec{x})$ as a subformula.

The reason why we can make such an assumption is that, given a quantifier-free
$\Lang$-formula $q(\vec{x},\vec{y})$, we may redefine $\bar{\chi}_{q}(\vec{x})$
as
\[
\bar{\chi}_{q}(\vec{x})=\bigwedge_{\vec{x}=(\vec{x}_{1},\vec{x}_{2})}\chi_{q(\vec{x}_{1},\cdot,\cdot)}(\vec{x}_{1}),
\]
where the conjunction ranges over all decompositions of the tuple
$\vec{x}$ into two tuples of variables (not necessarily in order).
\end{remark}

\subsection{Main Result}
\begin{thm}
\label{thm:main_final}If $D$ is a $0'$-computable c.e.-typed structure
over the signature $\Lang'\supseteq\Lang$ satisfying the QETP, then
the reduct $D\upto\Lang$ admits a computable copy $M$. There is
a $0'$-computable isomorphism between $D\upto\Lang$ and $M$. This
result is uniform in the index for $D$ and an index for witnesses
that $D$ satisfies the QETP,\footnote{This includes: An index for the function that takes a predicate $P$
in $\Lang$ to a predicate that is interpreted in $D$ as $\neg P$,
an index for the computable functions $\chi$ and $\etau$, and an
index for the $0'$-computable function $\tau$.} and relativizes uniformly.
\end{thm}
\begin{proof}
We build a computable copy of $D\upto\Lang$ via a finite injury argument.
First, some background assumptions:
\begin{itemize}
\item We see the structure $D$ as being given by a computable process,
which outputs elements and c.e.\ indices for the atomic $\Lang'$-type
of the elements output until now, and sometimes erases previously
constructed elements. We will be executing this construction in parallel
in the background of ours.
\item An element of $D$ which is added in this manner and never removed
will be referred to as a \emph{real element} of $D$.
\item We will build the $\Lang$-structure $M\cong D\upto\Lang$ via a computable
process that outputs a sequence of elements, and with each element,
the quantifier-free $\Lang$-type of the elements output until now.
\item We well-order the elements of $M$ and the elements of $D$ by time
of addition. This includes ``fake'' elements of $D$. When we refer
to ``the first element of $M$ (or $D$) that satisfies such-and-such
property'', we mean the one that was added earliest.
\item We will imagine an assortment of agents, referred to as ``agent $i$''
for $i=0,1,2,\dots$, each of which will be responsible for one step
in a back-and-forth construction of an isomorphism between $D$ and
the structure $M$ that we are constructing. Roughly speaking, agent
$i$ will be responsible for satisfying the requirement: The $i$-th
real element of $D$ is matched to an element of $M$, and the $i$-th
element of $M$ is matched to an element of $D$, and the matchings
created by agents $0$ to $i$ form an isomorphism of finite $\Lang$-structures.
We describe requirement $i$ in more detail below, by describing how
it may be injured, after we have introduced some further notation
about the construction.
\item We assume that only one agent is active at each time. When agent $i$
is done with its task, it will activate agent $i+1$. When an injury
occurs, we divert execution back to a prior agent, say agent $j$,
and delete the state of all agents past agent $j$.
\item The internal state of this construction, that is, the information
we will be keeping track of, consists of the following data (plus
whatever auxiliary information is necessary to keep track of the agent
mechanism):
\begin{itemize}
\item The current status of the structures $M$ and $D$,
\item One-to-one matchings between some elements of $M$ and some elements
of $D$, the latter of which may or may not be real,
\item To some elements $m$ of $M$, we match a triple of formulas $(q(\vec{x},y,\vec{z}),\varphi(\vec{x}),\tau(\vec{x},y))$
that have been output by the process from \rmkref{tauwork}, together
with a decomposition of the elements of $M$ at the time that this
matching was created into tuples $\vec{x}$, $y=m$, and $\vec{z}$.
\item To each matching of the previous two kinds, we keep track of which
agent created it and when.
\end{itemize}
\item The internal state of this construction may be modified via the following
operations:
\begin{itemize}
\item We may add a new element to $M$, in which case we immediately match
it to an element of $D$,
\item We may match a triple of formulas $(q,\varphi,\tau)$ to an unmatched
element $y\in M$, together with the decomposition $M_{\mathrm{current}}=(\vec{x},y,\vec{z})$,
where $\vec{x}$ denotes the elements of $M$ that have a match in
$D$ at this moment,
\item We may create a matching between an element of $M$ and an element
of $D$, if the element of $M$ has a triple $(q,\varphi,\tau)$ assigned,
\item In the event of injury, we may remove matches of the three kinds above.
\end{itemize}
\item We say requirement $i$ is injured if any of the following three things
happen:
\begin{itemize}
\item An element of $D$ which was matched by agent $i$ to an element of
$M$ is deleted,
\item There is an element of $M$, to whom agent $i$ matched a triple of
formulas $(q,\varphi,\tau)$, for which the triple of formulas has
been deleted from the enumeration of the process of \rmkref{tauwork},
\item There is an element of $M$, to whom agent $i$ matched a triple of
formulas $(q,\varphi,\tau)$, and a matched element $d\in D$, such
that our guess about whether $\tau$ holds of this element has changed
to `no'.
\end{itemize}
In any of these cases, all matches created after the ``faulty match''
will be deleted.
\end{itemize}
Now we list some conditions to be maintained, barring injury. Together
with each assumption, we point out the operations that we need to
be careful in exercising -- the remaining operations will always
preserve the given assumption. Moreover, the assumptions will not
be broken with the passage of time unless injury happens, in which
case the injury-handling part of the algorithm will rectify any issues.
\begin{enumerate}[label={Assumption (\Alph*)}, leftmargin=*]
\item \label{enu:asmchi}If $M=(\vec{x},\vec{y})$, where $\vec{x}$ denotes
the elements of $M$ that currently admit matches in $D$, $\vec{b}$
denotes these matches, and $q(\vec{x},\vec{y})$ is the type of the
elements of $M$, we demand that $D\vDash\chi_{q}(\vec{b})$.

We need to be careful to ensure that this is true when we add new
elements to $M$, and when we create matches between elements of $M$
and elements of $D$. We would also need to be careful in the event
that matches are deleted, but this is where \rmkref{weakerchi} becomes
relevant, as it ensures that this is not a concern.
\item \label{enu:asmphi}If $y\in M$ is matched with a triple of formulas
$(q,\varphi,\tau)$ and decomposition $(\vec{x},y,\vec{z})\subseteq M$,
we demand that $M\vDash q(\vec{x},y,\vec{z})$ and that all elements
of $\vec{x}$ admit matches $\vec{b}$ in $D$, which satisfy $D\vDash\varphi(\vec{b})$.

We need to be careful to ensure that this is true when we assign a
triple to an element. We don't need to be careful when removing matches
or elements of $D$, because our injury process will, in that case,
unmatch the formulas assigned to $y$.
\item \label{enu:asmutau}At most one unmatched element $m$ of $M$ will
have an assigned triple at any given time, and in this event, the
variables $\vec{x}$ correspond to the elements of $M$ currently
matched to elements of $D$, and the next action to be taken will
be to match $m$ to an element of $D$.
\item \label{enu:asmtau}If $y$ is an element of $M$ with assigned triple
$(q(\vec{x},y,\vec{z}),\varphi(\vec{x}),\tau(\vec{x},y))$ and $y$
admits a match $d$ in $D$, if $\vec{b}$ is the tuple of elements
matched with $\vec{x}$, it must be the case that $D\vDash\tau(\vec{b},d)$
\emph{at the current approximation} (in the sense of \rmkref{llorc2}).

We need to be careful to ensure that this is true when we create a
match between an element of $M$ with an assigned triple $(q,\varphi,\tau)$
and an element of $D$. We also need to be careful if the current
approximation (in the sense of \rmkref{llorc2}) ever changes, which
is why `the approximation of the value of $\tau(\vec{b},d)$ changes'
is a cause of injury.
\item \label{enu:asmetau}If $m$ is an element with assigned formulas $(q(\vec{x},y,\vec{z}),\varphi(\vec{x}),\tau(\vec{x},y))$,
with $\vec{b}$ denoting the matches of the $\vec{x}$, and at this
stage $M$ is decomposed as $M=(\vec{x},y,\vec{z},\vec{w})$ with
quantifier-free $\Lang$-type $Q(\vec{x},y,\vec{z},\vec{w})$, it
must be the case that, if the $\vec{b}$ are real elements of $D$,
we have $D\vDash\etau_{q\varphi Q}(\vec{b})$.

We need to be careful to ensure that this is true when new elements
are added to $M$, and when we assign triples $(q,\varphi,\tau)$
to elements of $M$.
\end{enumerate}
The process starts with the initialization of agent $0$.

Now we describe the algorithm that will be followed by agent $i$.
When agent $i$ is initialized:
\begin{enumerate}[label={Step \Roman*.},ref={Step \Roman*}]
\item \label{enu:mainstep1}Find the first element of $M$, if any, that
does not have a match in $D$. If there is none, skip to \ref{enu:mainstep2}.
If there is one, call it $m$.
\item If $m$ has an assigned triple of formulas $(q,\varphi,\tau)$, skip
to \ref{enu:mainstep1.4}.\footnote{Technically, this instruction is unnecessary. The only way that the
condition could be met is in the case of injury (namely, if the triple
$(q,\varphi,\tau)$ assigned to $m$ is removed from the enumeration
from \rmkref{tauwork}), but in the description of injury-handling
we already skip straight to \ref{enu:mainstep1.4} regardless.}
\item \label{enu:mainstep1.3}If $m$ does not have an assigned triple:
Let $(\vec{x},y,\vec{z})$ be the current elements of $M$, with $\vec{x}$
corresponding to the elements that currently have a match in $D$
-- call their matches $\vec{b}$ --, $y$ corresponding to $m$,
and $\vec{z}$ corresponding to the remaining elements. Look through
the triples $(q,\varphi,\tau)$ enumerated as per \rmkref{tauwork}
until you find such a triple that satisfies $M\vDash q(\vec{x},y,\vec{z})$
and $D\vDash\varphi(\vec{b})$. If all elements of $\vec{b}$ are
real, such a triple will be found in finite time by \enuref{asmchi}
in conjunction with Item \ref{enu:qetp2} of the QETP. If not all
are real, injury will happen. Once such a triple is found, we assign
to $m$ the triple $(q,\varphi,\tau)$ with the decomposition $(\vec{x},y,\vec{z})$
above.

We need to make sure that \ref{enu:asmphi} and \ref{enu:asmetau}
are preserved.
\begin{itemize}
\item \ref{enu:asmphi} -- By design.
\item \ref{enu:asmetau} -- We need to guarantee that, if the $\vec{b}$
are all real, we have $D\vDash\etau_{q\varphi q}(\vec{b})$. This
is a consequence of the note at the end of Item \ref{enu:qetpetau}
of the QETP.
\end{itemize}
\item \label{enu:mainstep1.4}Let $(q,\varphi,\tau)$ be the formula assigned
to $m$. The formula $q(\vec{x},y,\vec{z})$ will not, in general,
be the quantifier-free $\Lang$-type of the current elements of $M$,
but it is definitely the quantifier-free $\Lang$-type of some subset
of elements of $M$. By \ref{enu:asmutau}, the elements corresponding
to $\vec{x}$ are exactly the ones that currently admit matches in
$D$. Let $\vec{b}$ denote their matches.

We let $Q(\vec{x},y,\vec{z},\vec{w})$ denote the quantifier-free
type of the current elements of $M$, where $\vec{x}$, $y$, and
$\vec{z}$ are the same as above.

Look through the elements of $D$ until you find an element $d$ satisfying
$D\vDash\tau(\vec{b},d)$ to current approximation (in the sense of
\rmkref{llorc2}).

{\small Technical note: This search must be done in some systematic
way so that, if all $\vec{b}$ are real, we eventually find a }{\small\emph{real}}{\small{}
element $d$ satisfying this formula in such a way that the approximation
stays fixed. An example of such a systematic way: Let's say we give
the $n$-th element added to $D$, say $d_{n}$, a penalty of $n$
stages in the time it takes to compute whether $D\vDash\tau(\vec{b},d_{n})$,
and select the $d$ such that $D\vDash\tau(\vec{b},d)$ is verified
(and not disproven until now) in the least amount of stages. This
will be important when proving finite injury.}{\small\par}

We claim that, unless an injury occurs, this search will eventually
terminate. By this we mean: Assume that all $\vec{b}$ are real, and
that the triple $(q,\varphi,\tau)$ will never be removed from the
enumeration from \rmkref{tauwork}. Then, by \enuref{asmchi} it must
be the case that $D\vDash\chi_{Q}(\vec{b})$ and by \enuref{asmetau}
it must be the case that $D\vDash\etau_{q\varphi Q}(\vec{b})$. Thus,
by Item \ref{enu:qetpetau} of the QETP, it must be the case that
$D\vDash(\exists y)\tau_{q}(\vec{b},y)$, and any witness to this
statement will serve as the $d$ above.

Once such a $d$ is found, we match $m$ to $d$.

We need to make sure that \enuref{asmchi} and \enuref{asmtau} hold.
\begin{itemize}
\item \enuref{asmchi} -- If $d$ is real, this is direct by Item \ref{enu:qetp3}
of the QETP.
\item \enuref{asmtau} -- By construction
\end{itemize}
\item \label{enu:mainstep2}Next, find the first element of $D$ -- call
it $d$ -- that does not have a match in $M$. If there is no such
element as of right now, proceed to \ref{enu:laststep}.

Let $q(\vec{x},\vec{y})$ be the quantifier-free $\Lang$-type of
the current elements of $M$, with $\vec{x}$ denoting the elements
that currently have matches in $D$, and let $\vec{b}$ denote their
matches. Let $r(\vec{x},z)$ be the quantifier-free $\Lang$-type
of the tuple $(\vec{b},d)$.\footnote{Here, we use Item \ref{enu:qetp0} of the QETP to computably obtain
(a strong index for) $r(\vec{x},z)$ from a c.e.\ index for the atomic
$\Lang'$-type of $(\vec{b},d)$. More generally, we use the fact
that Item \ref{enu:qetp0} of the QETP allows us to computably obtain
a strong index for the $\Lang$-type of a tuple from a c.e.\ index
for its $\Lang'$-type.} In parallel, look through the quantifier-free $\Lang$ existential
types $Q(\vec{x},z,\vec{y})$ until you find such a $Q$ extending
$q(\vec{x},\vec{y})\land r(\vec{x},z)$ such that:
\begin{itemize}
\item $D\vDash\chi_{Q(\vec{x},z,\cdot)}(\vec{b},d)$, and
\item For every $m$ with an assigned triple $(q_{0}(\vec{x}_{0},y_{0},\vec{z}_{0}),\varphi(\vec{x}_{0}),\tau(\vec{x}_{0},y_{0}))$,
if $\vec{b}_{0}$ are the elements of $M$ corresponding to $\vec{x}_{0}$,
$D\vDash\etau_{q_{0}\varphi Q}(\vec{b}_{0})$.
\end{itemize}
We argue that either such a $Q$ is found in finite time or otherwise
an injury will occur. Indeed, suppose that both $d$ and the elements
of $\vec{b}$ are all real, and that every matched triple $(q_{0},\varphi,\tau)$
will never be injured. Now, since we assume $D\vDash\chi_{q}(\vec{b})$,
there must be a tuple of real elements $\vec{c}$ in $D$ such that
$D\vDash q(\vec{b},\vec{c})$. Now, we claim that the type $Q$ of
the tuple $(\vec{b},d,\vec{c})$ satisfies the desired requisites:
\begin{itemize}
\item Since $D\vDash Q(\vec{b},d,\vec{c})$, we must have $D\vDash(\exists\vec{y})Q(\vec{b},d,\vec{y})$
and hence $D\vDash\chi_{Q}(\vec{b},d)$, and
\item If $m$ has the triple $(q_{0},\varphi,\tau)$ assigned, with decomposition
$(\vec{x}_{0},y_{0},\vec{z}_{0})\subseteq M$, with $\vec{b}_{0}$
being the elements of $D$ corresponding to $\vec{x}_{0}$ and $d_{0}$
being the element corresponding to $m$, we have $D\vDash\chi_{Q}(\vec{b}_{0},d_{0})$
as a consequence of $D\vDash Q(\vec{b},d,\vec{c})$, and moreover
$D\vDash\tau(\vec{b}_{0},d_{0})$ by \enuref{asmtau}, hence $D\vDash(\exists y)\left[\chi_{Q}(\vec{b}_{0},y)\land\tau_{q\varphi}(\vec{b}_{0},y)\right]$,
and thus $D\vDash\etau_{q_{0}\varphi Q}(\vec{b}_{0})$ by Item \ref{enu:qetpetau}
of the QETP.
\end{itemize}
Once such a $Q(\vec{x},z,\vec{y})$ is found, one of two things occurs.
Either it includes $z=y_{j}$ for some value of $j$, in which case
we match $d$ to the element corresponding to $y_{j}$, or it includes
the information that all its variables represent distinct elements,
in which case we add a new element to $M$, whose relations to the
previous elements are those dictated by $Q$, and match it to $d$.

We need to make sure that \enuref{asmchi} and \enuref{asmetau} hold.
\begin{itemize}
\item \enuref{asmchi} -- By construction, since we ensured $D\vDash\chi_{Q}(\vec{b},d)$.
\item \enuref{asmetau} -- By construction.
\end{itemize}
\item \label{enu:laststep}Once all the above steps are done, agent $i$
will pause its execution, and initialize agent $i+1$.
\end{enumerate}
We now describe the process of injury. There are three means of injury:
\begin{itemize}
\item Suppose that an element of $D$ is seen to be removed. If this element
does not have a match in $M$, no action is taken. On the other hand,
suppose that this element $d$ is matched to an element $m$ in $M$.
In that event, suppose that agent $i$ is the one that made this match,
and halt the execution of all agents with index greater than $i$,
erasing matches and assigned formulas created past the moment that
$m$ and $d$ were matched. If the match between $d$ and $m$ was
created by \ref{enu:mainstep1}, restart the execution of agent $i$.
If the match between $d$ and $m$ was created by the \ref{enu:mainstep2},
initialize agent $i+1$ (and do not restart agent $i$).
\item Suppose that an element $m\in M$ has an assigned triple $(q,\varphi,\tau)$,
which is seen to be deleted from the enumeration described in \rmkref{tauwork}
at a given stage. Suppose agent $i$ is the one who has assigned this
formula. Then, we halt the execution of all agents with index greater
than $i$, and erase all matches and assigned formulas created past
the moment that this formula was assigned to $m$. Then, execute agent
$i$ starting in \ref{enu:mainstep1.3}.
\item Suppose that an element $m\in M$ has an assigned triple $(q(\vec{x},y,\vec{z}),\varphi(\vec{x}),\tau(\vec{x},y))$
and a matched element $d$. Suppose agent $i$ is the one who has
created these matches. Suppose that the elements $\vec{x}$ are matched
to the tuple $\vec{b}$, and suppose that the current approximation
of $D\vDash\tau(\vec{b},d)$ (in the sense of \rmkref{llorc2}) has
changed to `no'. Then, we halt the execution of all agents with
index greater than $i$, erase all matches and assigned formulas created
past this point (including the match between $m$ and $d$, but we
retain the assigned triple $(q,\varphi,\tau)$), and execute agent
$i$ starting in \ref{enu:mainstep1.4}.
\end{itemize}
This concludes the construction. It remains to show that the resulting
structure $M$ is isomorphic to $D\upto\Lang$, and this is done via
a finite injury argument. We show that each requirement will be injured
finitely many times. This will show that every element in $D$ is
eventually matched to an element of $M$ (the $i$-th element of $D$
(counting both real and fake elements) will be matched by, at worst,
the $i$-th agent past the moment when this element is added to $D$),
and that every element of $M$ is eventually matched to an element
of $D$ (given an orphan element of $M$, inductively consider the
first agent activated after all previous elements of $M$ have been
matched), and since these matches all preserve the predicates in $\Lang$,
the resulting map will be an isomorphism.

Suppose, for the sake of induction, that all requirements prior to
$i$ are injured finitely many times. We show that requirement $i$
itself also suffers finite injury. We do this individually for each
of the matches potentially created by agent $i$:
\begin{itemize}
\item First, we show that the matching created by the first orphan element
of $m$ at this moment, if there is any, is injured finitely many
times. Indeed, per the notation of \rmkref{tauwork} and using Item
\ref{enu:qetp2} of the QETP, eventually a triple $(q,\varphi,\tau)$
will be enumerated and never removed, such that the relevant elements
of $D$ satisfy $\chi_{q}\land\varphi$. Past this point, after all
previously-added false triples have been removed, this triple (or
another) will be permanently assigned to $m$. Afterward, we know
(see the technical note in \ref{enu:mainstep1.4}) that after enough
steps we match $m$ with a real element $d$ of $D$ that satisfies
$\tau(\vec{b},d)$ at the current approximation, which will then stay
fixed. This match will never cause injury.
\item Now, assuming that we've reached a time where the matching discussed
in the previous bullet point or created by previous agents will never
be removed, suppose that there is an unmatched element in $D$ for
agent $i$ to match with an element of $M$. If this element is real,
there is no injury. If the element is not real, then it is eventually
erased. Then execution will move to agent $i+1$, and requirement
$i$ is never injured.
\end{itemize}
This concludes the proof.
\end{proof}

\section{\protect\label{sec:prevs_applications}Previously Known Applications}

In this section, we list some known results that can be obtained using
\thmref{main_final}. We give the proof just for the most interesting
one.

\subsection{Elementary results}
\begin{thm}[{{\cites[VII.26]{montalban_cst}}}]
\label{thm:JIlinear-1}Suppose that $L$ is a linear order for which
every element admits a successor and predecessor, and suppose that
the structure $(L,<,S)$ admits a $0'$-computable copy, where $S(x,y)$
is the successor relation. Then, the structure $(L,<)$ admits a computable
copy.
\end{thm}
\begin{proof}
Omitted. See \cite{self_escaping_tennenbaum}.
\end{proof}
\begin{prop}[{Folklore, {\cites[Theorem 9.1.4]{cst_downey_melnikov}}}]
\label{prop:eqvinf}Let $(E,\eqvrel)$ be an equivalence relation
with infinitely many infinite equivalence classes, and assume that
the structure $(E,\eqvrel,\{P_{\geq n}\}_{n\in\N})$ admits a $0'$-computable
copy, where $P_{\geq n}(x)$ is the predicate ``the equivalence class
of $x$ has size at least $n$''.\footnote{This includes the case where the equivalence class of $x$ is infinite.}
Then, $(E,\eqvrel)$ admits a computable copy.
\end{prop}
\begin{proof}
Omitted. See \cite{self_escaping_tennenbaum}.
\end{proof}
\begin{definition}[{Khisamiev, {\cites[Definition 9.1.3]{cst_downey_melnikov}}}]
\label{def:limitwisemonotonic}A set $X\subseteq\N$ is said to be
\emph{limitwise monotonic} if there is a total computable function
$f\colon\N\times\N\to\N$ such that, for every $x\in\N$, the sequence
$\{f(x,y)\}_{y\in\N}$ is weakly increasing (i.e.\ nondecreasing)
and bounded (and therefore convergent), and $X$ is the image of the
pointwise limit function $g(x)=\lim_{y}f(x,y)$.
\end{definition}
%
\begin{definition}[{{\cites[Definition 9.1.1]{cst_downey_melnikov}}}]
Let $(E,\eqvrel)$ be an equivalence relation. The \emph{characteristic
set of $E$} is the set
\[
\#E=\{\,n\in\N\mid\text{\ensuremath{E} admits at least one equivalence class of size \ensuremath{n}}\,\}.
\]
\end{definition}
\begin{prop}[{Folklore, {\cites[Theorem 9.1.4]{cst_downey_melnikov}}}]
\label{prop:eqvfin}Let $(E,\eqvrel)$ be an equivalence relation
with finitely many infinite equivalence classes, and assume that the
structure $(E,\eqvrel,\{P_{\geq n}\}_{n\in\N})$ admits a $0'$-computable
copy, where $P_{\geq n}(x)$ is the predicate ``the equivalence class
of $x$ has size at least $n$''. Assume moreover that the characteristic
set of $E$ is limitwise monotonic. Then, $(E,\eqvrel)$ admits a
computable copy.
\end{prop}
\begin{proof}
Omitted. See \cite{self_escaping_tennenbaum}.
\end{proof}
\begin{prop}[{{\cites[Proposition 3.8]{strong_jump_inversion}}}]
\label{prop:trees}
\global\long\def\leaf{\mathop{\mathrm{leaf}}}%
Let $T$ be an infinite subtree of $\omega^{<\omega}$ such that every
node of $T$ that is not a leaf has infinitely many successors, of
which finitely many are leaves. For the signature, we use the successor
relation $S(x,y)$ and a unary relation $R_{0}(x)$ for the root.
The other important predicate for a structural jump is $\leaf(x)$
($x$ is a leaf).

If the structure $(T,\leaf,S,R_{0})$ admits a $0'$-computable copy,
then the structure $(T,S,R_{0})$ admits a computable copy.

\end{prop}
\begin{proof}
Omitted. See \cite{self_escaping_tennenbaum}.
\end{proof}

\subsection{Khisamiev's Theorem}

The following theorem is in the literature, see e.g. Chapter 5 of
\cite{cst_downey_melnikov}.

In the following, we interpret a c.e.\ presentation of a group as
being a description of the group as a list of generators and relations,
such that there is a computer program that enumerates the list of
generators and relations. In the present case, this is equivalent
to the more general notion of c.e.\ presentation of a structure in
computability theory.
\begin{thm}[Khisamiev]
\label{thm:khisamiev}Every c.e.\ presented torsion-free abelian
group is isomorphic to a computable group. Furthermore, if the group
is non-trivial, then this computable copy can be built uniformly in
the index of the c.e.\ presentation.
\end{thm}
\begin{proof}
Given a c.e.\ presentation of a group $G$, we first build a $0'$-computable
c.e.-typed copy of $G$ over an appropriate signature $\Lang'$, and
we will then show that any nontrivial torsion-free abelian group has
the QETP (uniformly) over $\Lang\subseteq\Lang'$, where $\Lang$
is the signature containing only the ternary predicate $S(x,y,z)\equiv\text{``\ensuremath{x+y=z}''}$(and
its negation, and equality and its negation).

The signature $\Lang'$ contains $\Lang$ plus:
\begin{itemize}
\item The countably many predicates $Q_{n,\vec{c}}(\vec{x})$, parametrized
over $n\in\N$ and finite tuples $c_{1},\dots,c_{k}\in\Z$:
\[
Q_{n,\vec{c}}(\vec{x})\equiv\text{``\ensuremath{n\mid c_{1}x_{1}+\dots+c_{k}x_{k}}''}.
\]
Note that we allow $n$ to equal zero, in which case (since $0\mid y$
iff $y=0$) $Q_{0,\vec{c}}(\vec{x})$ corresponds to the predicate
$\vec{c}\cdot\vec{x}=0$.
\item The countably many predicates $\neg Q_{0,\vec{c}}(\vec{x})$, parametrized
over finite tuples $c_{1},\dots,c_{k}\in\Z$:
\[
\neg Q_{0,\vec{c}}(\vec{x})\equiv\text{``\ensuremath{c_{1}x_{1}+\dots+c_{k}x_{k}\neq0}''}.
\]
\item Countably many predicates $\{n\Delta(x)\}_{n\in\Z}$, whose meaning
is as follows: We will carefully pick a single designated nonzero
element of $G$, which we refer to as $\delta$, and $n\Delta(x)$
should be interpreted as ``$x=n\delta$''.
\end{itemize}
We need some linear algebra of $\Z$-modules to show that, from a
c.e.\ presentation of $G$, we can uniformly obtain a $0'$-computable
c.e.-typed copy of $G$. Then, we will explain how we choose the element
of $G$ that will be tagged with $\Delta$. To obtain the $0'$-computable
c.e.-typed copy of $G$, here is the basic idea: Consider the computable
process that outputs all possible linear combinations of generators.
We want to, with each new linear combination of generators, say $d$,
first check whether it is equal to any of the combinations previously
produced (this is an easy check with $0'$), and if it is not, we
want to output a c.e.\ index for its type relative to the elements
previously output. This is easy for the positive predicates -- to
check if, for example, $5\mid4d+6c$ (where $c$ is a previously output
combination of generators), it suffices to check if there is a further
combination of generators $u$ that satisfies the relation $5u-4d-6c=0$.
This is a c.e.\ check. The difficulty lies in the negative predicates,
of which $\neg Q_{0,\vec{c}}(\vec{x})$ is the general case. In other
words, we need to encode \emph{into a single c.e.\ index} the full
information about what linear combinations of elements-output-so-far
equal zero.

Here is a fact from commutative algebra, whose proof can be found
in any standard book that covers modules over principal ideal domains
(e.g.\ \cite{dummit_foote}):
\begin{enumerate}[label={Fact \Alph*:},ref={Fact \Alph*}]
\item Any submodule $M$ of $\Z^{k}$ admits a finite basis.
\end{enumerate}
Here is the effective version of this fact that we will use:
\begin{enumerate}[resume, resume*]
\item \label{enu:cemodulebasis}Given a c.e.\ index for a submodule $M$
of $\Z^{k}$, we may $0'$-effectively find a basis of $M$.
\end{enumerate}
Proof idea: Consider the $k$-width and possibly-infinite-height matrix
whose rows are the elements of $M$. Use $0'$ to find the GCD of
the elements of the first column, call it $d$. If $d=0$ (i.e.\ all
elements of the first column are zero), skip to the next step. If
$d\neq0$, find a combination of rows of $M$ that has $d$ in its
first coordinate, and add this combination to the basis. Then, if
$k\geq1$, recursively apply this algorithm to the $(k-1)$-width
and possibly-infinite-height matrix whose elements consist of the
rows of the previous matrix that have a $0$ as a first entry. Upon
finding a combination of rows that has the new GCD as its first coordinates,
add the corresponding combination of the rows of the original matrix
to the basis. After $k$ steps, we have a basis of $M$.

Returning to our original problem: We have a collection of combinations
of generators, say $d_{0}$, $d_{1}$, $\dots$, $d_{k}$, and we
wish to encode all possible linear combinations of these elements
that are equal to zero. We do this by applying \ref{enu:cemodulebasis}
to the c.e.\ submodule of $\Z^{k+1}$ given by
\[
M=\{\,\vec{c}\in\Z^{k+1}\mid\vec{c}\cdot\vec{d}=0\,\}.
\]
With access to a basis of $M$, say $b_{1},\dots,b_{\ell}$ (encoded
into the columns of a matrix $B$), we can tell whether any set of
coefficients $\vec{c}$ is in $M$ (and hence, whether we should say
yes or no to $Q_{0,\vec{c}}(\vec{d})$) by using standard algorithms
to solve the linear equation $B\vec{x}=\vec{c}$ over $\Q$, and conclude
$\vec{c}\in M$ if and only if a solution exists and its coefficients
are all integers. Coding this basis into a c.e.\ index, we obtain
a c.e.\ index for the atomic $\Lang'$-type of a tuple $\vec{d}$
of linear combinations of generators.

\smallskip{}

Now, let us explain how to choose the element $\delta$, which we
will use to decide the predicates $n\Delta$. Of course, this has
to have been done in advance, so that we can hardcode the combination
of generators corresponding to $\delta$ into the c.e.\ indices of
the types of the elements of $G$ in advance.

There is an easy way and a hard way to choose $\delta$. The easy
way is simply to pick an arbitrary nonzero element of $G$. The issue
is that to make this choice requires the usage of an oracle for $0'$,
which will lead to a non-uniform presentation. Here is a uniform way
to choose the element. Let $\{C_{i}\}_{i\in\N}$ enumerate all possible
combinations of generators. To start, we will guess $\delta=C_{0}$,
and look through the relations until (if ever) we see enough relations
to conclude $C_{0}=0$. When this happens, we will guess $\delta=N_{1}C_{1}$,
where $N_{1}$ is a nonzero number divisible by enough integers. This
choice is made so that, for every number $k$ for which we'd been
able to conclude $k\mid C_{0}$ before we found that $C_{0}=0$, we
have $k\mid N_{1}$ and so $k\mid N_{1}C_{1}$. We proceed in this
manner, enumerating relations until we are able to conclude that $N_{1}C_{1}=0$,
or equivalently that $C_{1}=0$. If this happens, we change our mind
to guess $\delta=N_{2}C_{2}$, where $N_{2}$ has enough divisibilities
that, again, every divisibility we'd already concluded about $N_{1}C_{1}$
also holds for $N_{2}$ and hence $N_{2}C_{2}$. Since, by assumption,
$G$ is a nontrivial group, eventually this procedure will run into
a nonzero combination of generators, and $0'$ is able to figure out
which element is $\delta$ in advance. The important fact about this
way of choosing $\delta$ is that we can \emph{uniformly} find a c.e.\ index
for the type of $\delta$.

\medskip{}
Now that we have a $0'$-computable c.e.-typed copy of our group over
the signature $\Lang'$, we will show that every nontrivial torsion-free
abelian group with a designated nonzero element with c.e.\ atomic
type satisfies the QETP uniformly.
\begin{itemize}
\item First, we construct $\chi_{q}$. This is an adaptation of the standard
proof that the theory of torsion-free abelian groups with divisibility-by-integer
predicates admits quantifier elimination, which can be found e.g.\ in
\cites[Appendix A.2]{hodges_model_theory}. Very briefly:
\begin{lemma}
Any existential statement is equivalent to a system of the type
\begin{equation}
(\exists\vec{y})q(\vec{x},\vec{y})\Leftrightarrow\left\{ \begin{array}{l}
M_{0}\vec{x}=0,\\
n_{1}\mid\vec{c}_{1}\cdot\vec{x},\\
\vdots\\
n_{\ell}\mid\vec{c}_{\ell}\cdot\vec{x},\\
K_{0}\vec{x}\neq0,
\end{array}\right.\label{eq:syseqchi}
\end{equation}
with the system found effectively. We define $\chi_{q}(\vec{x})$
to be the right-hand side of (\ref{eq:syseqchi}).
\end{lemma}
The effective transformation goes as follows:
\begin{itemize}
\item Express the type $q(\vec{x},\vec{y})$ as a system of integer coefficient
equations and inequations
\[
q(\vec{x},\vec{y})\equiv\left\{ \begin{array}{l}
M\vec{x}=N\vec{y},\\
K\vec{x}\neq L\vec{y}.
\end{array}\right.
\]

\textbf{Important:} The notation $K\vec{x}\neq L\vec{y}$ means ``every
entry of the vector $K\vec{x}$ is distinct from the corresponding
entry in $L\vec{y}$''.
\item Now, write $N=ADB$, where $D$ is a diagonal matrix padded with zeros
(i.e.\ of the form $\left[\begin{smallmatrix}D_{0} & 0\\
0 & 0
\end{smallmatrix}\right]$ with $D_{0}$ diagonal) and $A$ and $B$ are invertible, all with
integer entries \cites[Theorem 2.12]{bapat_graphs_matrices}. Then,
$q(\vec{x},\vec{y})$ is logically equivalent to a new system of equations
and inequations
\[
q(\vec{x},\vec{y})\Leftrightarrow\left\{ \begin{array}{l}
M'\vec{x}=DB\vec{y},\\
K'\vec{x}\neq LB^{-1}B\vec{y}.
\end{array}\right.
\]
\item The existence of a solution to this system is, since $B$ is invertible,
equivalent to the existence of a solution to the following other system
\[
(\exists\vec{y})q(\vec{x},\vec{y})\Leftrightarrow(\exists\vec{z})\left\{ \begin{array}{l}
M'\vec{x}=D\vec{z},\\
K'\vec{x}\neq L'\vec{z}.
\end{array}\right.
\]
\item The rows of the equation $M'\vec{x}=D\vec{z}$ that have nonzero entries
of $D$ yield equations of the type $nz_{i}=\vec{a}\cdot\vec{x}$.
This allows us to remove (while preserving logical equivalence) every
instance of $z_{i}$ from every other (in)equation: Any other (in)equation
can be multiplied by $n$ preserving logical equivalence (this uses
$n\neq0$ and the fact that we're working with torsion-free groups),
resulting in an instance of $\text{(coefficient)\ensuremath{\times}}nz_{i}$,
which can be substituted for $\text{(coefficient)}\times(\sum a_{i}x_{i})$.
Thus, we may eliminate all $z$ variables corresponding to nonzero
rows of $D$, and our system now becomes
\[
(\exists\vec{y})q(\vec{x},\vec{y})\Leftrightarrow(\exists\vec{z})(\exists\vec{w})\left\{ \begin{array}{l}
M_{0}\vec{x}=0,\\
n_{1}z_{1}=\vec{c}_{1}\cdot\vec{x},\\
\vdots\\
n_{\ell}z_{\ell}=\vec{c}_{\ell}\cdot\vec{x},\\
K'\vec{x}\neq L'\vec{w}.
\end{array}\right.
\]
\item The equations including the $z$ variables become divisibility statements,
i.e.\ $(\exists z_{i})nz_{i}=\vec{c}\cdot\vec{x}\Leftrightarrow n\mid\vec{c}\cdot\vec{x}\equiv Q_{n,\vec{c}}(\vec{x})$,
so we now need only deal with the inequations $K'\vec{x}\neq L'\vec{w}$.
\item Now, let us stratify the equations $K'\vec{x}\neq L'\vec{w}$ as follows:
First, consider the equations that use no $w$-variable. Then, consider
the ones that use $w_{1}$ and no other. Then, consider the ones that
use only $w_{1}$ and $w_{2}$, and so on. The equations that use
no $w$-variable are important -- keep them. These are a requirement
on $\vec{x}$ that we cannot get rid of. Let us refer to these equations
as $K_{0}\vec{x}\neq0$. On the other hand, at every new stage of
the strata, the resulting equations will doubtlessly be satisfied
by some $w_{i}$. For example, if the first stratum has $e$ equations,
this is forbidding at most $e$ possible values of $w_{1}$. But the
group is nontrivial and hence infinite, whereby there is at least
one valid choice for $w_{1}$. Then, in the next stratum, contingent
on this choice of $w_{1}$, there are once again at most (number of
equations) forbidden values for $w_{2}$. Simply choose $w_{2}$ to
not be one of those. Continuing in this manner, we see that we can
ensure that all equations that include a $w$ variable can be automatically
satisfied and are therefore adding no information. This concludes
the effective transformation.
\begin{remark}
\label{rmk:nonzeroxi}If it is known that not every $x_{i}$ equals
zero, one might in fact choose the $w_{i}$ variables to be linear
combinations of the $x_{i}$. This will be important later.
\end{remark}
\end{itemize}
\item Next, we construct $\tau$. Intuitively, given $q(\vec{x},y,\vec{z})$,
we consider all possible assignments of $y$ of the form $y=\frac{1}{a}\vec{c}\cdot\vec{x}$
(if such an element exists). We will prove below, when establishing
the QETP, that assuming not every element of $\vec{x}$ is null there
is such a combination that provably satisfies $\chi_{q}(\vec{x},\frac{1}{a}\vec{c}\cdot\vec{x})$.
We will choose one such combination, and set $\tau_{q\varphi}(\vec{x},y)\equiv(ay=\vec{c}\cdot\vec{x}$).
If every element of $\vec{x}$ is null, we set $\tau(\vec{x},y)\equiv N\Delta(y)$
where $N$ is large enough to encompass all divisibilities demanded
by the tuple of elements $\vec{z}$ (e.g.\ if $q(y,z)\equiv(z+z=y)$,
we would set $y=2\delta$).

Let's discuss this process more formally. Given $q(\vec{x},y,\vec{z})$:
\begin{itemize}
\item If $q$ includes the data $y=0$, set $\tau_{q\top}(\vec{x},y)\equiv(y=0)$.
Otherwise:
\item Set $\varphi_{0}(\vec{x})\equiv\land_{i}(x_{i}=0)$ and set $\tau_{q\varphi_{0}}(\vec{x},y)\equiv\varphi_{0}(\vec{x})\land N\Delta(y)$,
where $N$ is positive and large enough to ensure $(\forall d\neq0)\chi_{q}(\vec{0},Nd)$.
\item For every nonzero integer $a$ and tuple of integers $\vec{c}$, set
$\varphi_{a\vec{c}}(\vec{x})\equiv\text{``\ensuremath{\chi_{q}(\vec{x},\frac{1}{a}\vec{c}\cdot\vec{x})}''}$
and $\tau_{q\varphi_{a\vec{c}}}(\vec{x},y)\equiv(ay=\vec{c}\cdot\vec{x})$.
We now explain precisely what we mean by $\chi_{q}(\vec{x},\frac{1}{a}\vec{c}\cdot\vec{x})$.

Since $q$ is assumed to be a type, we have by definition that $\chi_{q}(\vec{x},y)$
is a system of equations as in (\ref{eq:syseqchi}). By multiplying
all equations by $a$, we obtain an equivalent system of equations
where every instance of $y$ is being multiplied by $a$. In other
words, we can rewrite $\chi_{q}(\vec{x},y)\Leftrightarrow\psi(\vec{x},ay)$,
and as a consequence the formula ``$\chi_{q}(\vec{x},\frac{1}{a}\vec{c}\cdot\vec{x})$''
may be taken to mean
\[
\text{``\ensuremath{\chi_{q}(\vec{x},\frac{1}{a}\vec{c}\cdot\vec{x})}''}\equiv(a\mid\vec{c}\cdot\vec{x})\land\psi(\vec{x},\vec{c}\cdot\vec{x}).
\]

\end{itemize}
\item Finally, we construct $\etau$. Given $q(\vec{x},y,\vec{z})$ and
$\varphi(\vec{x})$ as one of the above cases, as well as $Q(\vec{x},y,\vec{z},\vec{w})$
extending $q$, we wish to find a quantifier-free expression that
is equivalent to $(\exists y)\left[\chi_{Q}(\vec{x},y)\land\tau_{q\varphi}(\vec{x},y)\right]$.
Intuitively, the idea is that in either case our formula $\tau$ determines
$y$ exactly (as a function of $\vec{x}$), so we can just ``plug
that $y$ into $\chi_{Q}(\vec{x},y)$''. To do this more precisely,
we split into cases:
\begin{itemize}
\item If $\tau_{q\varphi}(\vec{x},y)\equiv(ay=\vec{c}\cdot\vec{x})$, we
apply the same ``multiply every equation by $a$'' trick to write
a formula that corresponds to ``$\left(a\mid\vec{c}\cdot\vec{x}\right)\land\chi_{Q}(\vec{x},\frac{1}{a}\vec{c}\cdot\vec{x})$''.
The resulting formula is the desired $\etau$.
\item If $\tau_{q\varphi}(\vec{x},y)\equiv\varphi_{0}(\vec{x})\land N\Delta(y)$,
then, intuitively, we know that for a tuple to satisfy $\tau$ it
must be that all $x_{i}$ are zero, and that $y=N\delta$. Therefore,
the existential quantifier may be eliminated, and we wish to determine
whether $\chi_{Q}(\vec{0},N\delta)$. This is where it is important
that we have a c.e.\ index for the type of $\delta$ -- and, if
we desire a uniform result, that we have this index without recourse
to any oracle. Knowing the c.e.\ index for the type of $\delta$,
we can just consult whether $\chi_{Q}(\vec{0},N\delta)$ holds --
more precisely, there is a computable sequence $\{B_{n}\}_{n\in\N}$
such that
\[
\left\{ \begin{aligned} & G\nvDash\chi_{Q}(\vec{0},N\delta) & \Rightarrow & B_{n}=\bot\text{ for all \ensuremath{n},}\\
 & G\vDash\chi_{Q}(\vec{0},N\delta) & \Rightarrow & B_{n}=\top\text{ for all but finitely many \ensuremath{n}.}
\end{aligned}
\right.
\]
In this scenario, we set 
\[
\etau_{q\varphi Q}(\vec{x})\equiv\llor_{n}B_{n}.
\]
\end{itemize}
\end{itemize}
Now we verify that the above definitions are witnesses to the QETP.
\begin{enumerate}[label=(\alph*)]
\item ($\Lang$ is finite and closed under negation) Obvious.
\item ($G\vDash\chi_{q}(\vec{x})\leftrightarrow(\exists\vec{y})q(\vec{x},\vec{y})$)
By construction.
\item There are two bullet points here. In reverse order:
\begin{itemize}
\item ($G\vDash\chi_{q}(\vec{x})\land\varphi(\vec{x})\rightarrow(\exists y)\tau_{q\varphi}(\vec{x},y)$)
Direct from construction.
\item (Every tuple $\vec{b}$ of elements of $G$ that satisfies $\chi_{q}(\vec{b})$
satisfies some $\varphi(\vec{b})$ with $(q,\varphi)$ in domain of
$\tau$)
\begin{itemize}
\item If $\vec{b}=\vec{0}$ this is obvious.
\item If some \textbf{$b_{i}\neq0$}: We wish to show that there are $a\in\Z_{\neq0}$
and $\vec{c}\subseteq\Z$ such that $a\mid\vec{c}\cdot\vec{b}$ and
$\chi_{q}(\vec{b},\frac{1}{a}\vec{c}\cdot\vec{b})$. To this effect,
we use the following related lemma, which obviously implies the desired
claim.
\begin{lemma}
If $q(\vec{x},\vec{y})$ is a quantifier-free type in the smaller
signature $\Lang$, and $\vec{b}$ is a tuple in $G$ such that $G\vDash(\exists\vec{y})q(\vec{b},\vec{y})$,
and some element of $\vec{b}$ is nonzero, then there exist integer
vectors $\vec{c}_{1},\dots,\vec{c}_{\ell}$ and nonzero integers $a_{1},\dots,a_{\ell}$
such that, for every $i$, $a_{i}\mid\vec{c}_{i}\cdot\vec{b}$, and
$G\vDash q(\vec{b},\frac{1}{a_{1}}\vec{c}_{1}\cdot\vec{b},\dots,\frac{1}{a_{\ell}}\vec{c}_{\ell}\cdot\vec{b})$.
\end{lemma}
The proof of this lemma uses a change of variables, similar to the
one used in the construction of $\chi_{q}$, followed by noticing
\rmkref{nonzeroxi}. The details are omitted.
\end{itemize}
\end{itemize}
\item ($G\vDash(\exists y)\left[\chi_{Q}(\vec{x},y)\land\tau_{q\varphi}(\vec{x},y)\right]\leftrightarrow\etau_{q\varphi Q}(\vec{x})$)
True by construction.
\item ($G\vDash\etau_{q\varphi Q}(\vec{x})\land\tau_{q\varphi}(\vec{x},y)\rightarrow\chi_{Q}(\vec{x},y)$)
Essentially, this is true because $\tau_{q\varphi}(\vec{x},y)$ \emph{always}
determines $y$ uniquely in terms of $\vec{x}$. As such, if it is
the case that there is some $y_{0}$ that satisfies $\chi_{Q}(\vec{x},y_{0})\land\tau_{q\varphi}(\vec{x},y_{0})$
(that is, $\etau_{q\varphi Q}(\vec{x})$ holds), and it is the case
that some specific $y$ satisfies $\tau_{q\varphi}(\vec{x},y)$, then
it must be the case that $y=y_{0}$ and therefore $\chi_{Q}(\vec{x},y)$.
\end{enumerate}

We are in condition to apply \thmref{main_final} to the $0'$-computable
c.e.-typed copy of $G$ (plus designated element). Since this copy
was obtained uniformly from the original presentation, the result
is uniform.
\end{proof}


\section{Escaping Tennenbaum's Theorem with More Truths}

In this section, we give a positive answer to the question posed by
Pakhomov in \cite{pakhomov}, about whether there exist theories definitionally
equivalent to ``$\mathsf{PA}$ + all $\Pi_{n}$ truths'' that admit
computable nonstandard models.

As described in the introduction, we will be working primarily with
the theory $\ZFniptc$, which consists of the axioms of $\mathsf{ZF}$,
with the removal of the axiom of infinity, and with the addition of
the axiom of transitive closure, which for our purposes is the axiom
stating that every set is in some level $V_{\alpha}$ of the von~Neumann
hierarchy.

\subsection{Generalizing Pakhomov's Construction}

This section is dedicated to the proof of \thmref{nestedpakhomov}:
\begin{thm}
\label{thm:nestedpakhomov}There is a nested sequence of consistent
c.e.\ theories $\ZFniptc\subseteq T^{0}\subseteq T^{1}\subseteq T^{2}\subseteq\cdots$
satisfying the following properties:
\begin{itemize}
\item Each $T^{n}$ is in the signature containing $\in$ and predicates
$S^{0}$, $\dots$, $S^{n}$, with each $S^{i}$ being an $(i+2)$-ary
predicate symbol,
\item All of these extensions are conservative, in the sense that they contain
no additional theorems in the predicate $\in$,
\item The predicates $\in$ and $S^{n}$ are definable in terms of the other
within $T^{n}$, and
\item Given an $X'$-computable model $D$ of $T^{n}$ restricted to the
signature containing $S^{i}$, $\dots$, and $S^{n}$, there is an
$X$-computable copy $M$ of $D\upto(S^{i+1},\dots,S^{n})$.
\end{itemize}
\end{thm}
%
We start the proof by defining the sequence $\{T_{n}\}_{n\in\N}$
inductively. For the base case, we set $T^{0}$ to be $\ZFniptc$,
plus the axiom $(\forall x,y)[S^{0}(x,y)\leftrightarrow(x\notin y)]$.
It trivially (and, in the case of the last bullet point, almost vacuously)
satisfies all required conditions. Now, we suppose that $T^{n}$ has
already been defined, and set out to define $T^{n+1}$.

Within $T^{n}$, we define the ordinal-indexed sequence of predicates
$S_{\alpha}^{n+1}(\vec{x},y)$, where $\vec{x}$ is an $(n+1)$-tuple
of variables and $\alpha$ is an ordinal. We define $S_{\alpha}^{n+1}$
inductively as a compatible sequence of relations on $V_{\alpha}$
with $S_{\lambda}^{n+1}$ having an obvious definition for limit ordinals
$\lambda$, so it suffices to describe the successor step. To assist
us in that regard, we define the following sequence of class functions:

Given sets $A$, $mA$, an ordinal $\alpha$, and sets $S^{i}A\subseteq(A\cup\{\alpha\})^{i+2}$
for $i=0,\dots,n+1$,\footnote{Here, the expression $S^{i}A$ is merely notation (where $i$ is a
finite ordinal and $A$ is a set), intended to evoke the idea that
$S^{i}A$ is defining the relation $S^{i}$ between the elements of
$A$ and a new element. Note that $S^{0}A$ is not the same as $A$.} define an element $w^{n}(\alpha,mA,A,SA,\dots,S^{n}A)$ inductively
in $n$ as follows.
\begin{itemize}
\item $w^{0}(\alpha,mA,A)=(\alpha,(A,mA))$, and
\item $w^{n+1}(\alpha,mA,A,S^{0}A,S^{1}A,\dots,S^{n}A,S^{n+1}A)=w^{n}(\alpha,(S^{n+1}A,mA),A,S^{0}A,\dots,S^{n}A)$.
%
\end{itemize}
The main defining features of this sequence of class functions are:
\begin{itemize}
\item For every $n$, $(\alpha,mA,A,S^{0}A,\dots,S^{n}A)\mapsto w^{n}(\alpha,mA,A,S^{0}A,\dots,S^{n}A)$
is an injective definable class function, and
\item $w^{n}(\alpha,mA,A,S^{0}A,\dots,S^{n}A)$ is \emph{not} in $V_{\alpha}$.
\end{itemize}
Let us go back to defining $S_{\alpha}^{n+1}$ inductively, assuming
that we have already defined $S^{0}$ up to $S^{n}$, for $n\geq0$.
Assuming that $S_{\alpha}^{n+1}$ is already defined on tuples of
elements of $V_{\alpha}$, we define a ``good sequence'' $(\alpha,mA,A,S^{0}A,\dots,S^{n+1}A)$
as one that satisfies the following properties:
\begin{itemize}
\item $A$ is a subset of $V_{\alpha}$,
\item For $i=0,\dots,n+1$, the set $S^{i}A$ is a subset of $(A\cup\{\alpha\})^{i+2}$,
\item For $i=0,\dots,n$, the set $S^{i}A\cap A^{i+2}$ coincides with the
relation $S^{i}$ on $A$,
\item The set $S^{n+1}A\cap A^{n+3}$ coincides with the relation $S_{\alpha}^{n+1}$
on $A$,
\item The set $S^{0}A$ contains all pairs $(a,\alpha)$, $(\alpha,a)$,
and $(\alpha,\alpha$), with $a\in A$.
\item For $i=0,\dots,n$, the following condition holds: For every $(i+3)$-tuple
$(\vec{a},b)$ of elements of $A\cup\{\alpha\}$, if $\vec{a}\notin S^{i}A$,
then $(\vec{a},b)\in S^{i+1}A$.
\end{itemize}
We are now ready to define $S_{\alpha+1}^{n+1}$. Given an $(n+3)$-tuple
$\vec{x}$ of elements in $V_{\alpha+1}$,
\begin{itemize}
\item If all elements of the tuple are in $V_{\alpha}$, set $S_{\alpha+1}^{n+1}(\vec{x})\equiv S_{\alpha}^{n+1}(\vec{x})$,
\item If there is a good sequence $(\beta,mA,A,S^{0}A,\dots,S^{n+1}A)$,
for some $\beta\leq\alpha$, such that the element $w=w^{n+1}(\beta,mA,A,S^{0}A,\dots,S^{n+1}A)$
is in $V_{\alpha+1}$, and every entry of $\vec{x}$ is an element
of $A\cup\{w\}$, set $S_{\alpha+1}^{n+1}(\vec{x})$ to hold true
if, and only if, upon replacing every instance of $w$ in the tuple
$\vec{x}$ by $\beta$, the resulting tuple is in $S^{n+1}A$.

In the sequence, when such a choice of good sequence is clear from
context, denote by $\vec{x}^{*}$ the operation of replacing every
entry of $\vec{x}$ that equals $\beta$ by $w$, and let $\vec{x}_{*}$
denote the inverse operation. Thus, the definition of $S_{\alpha+1}^{n+1}$
in this case can be reworded as: $S_{\alpha+1}^{n+1}(\vec{x})\equiv[\vec{x}^{*}\in S^{n+1}A]$.

Remark: Note that, by definition, if $(\beta,mA,A,S^{0}A,\dots,S^{n+1}A)$
is a good sequence and $n\geq0$, we also have that $(\beta,(S^{n+1}A,mA),A,S^{0}A,\dots,S^{n}A)$
is also a good sequence (for a smaller value of $n$), and thus the
element $w=w^{n+1}(\beta,mA,A,S^{0}A,\dots,S^{n+1})=w^{n}(\beta,(S^{n+1}A,mA),A,S^{0}A,\dots,S^{n}A)$
will also satisfy: $S^{n}(\vec{x})\equiv[\vec{x}^{*}\in S^{n}A]$
for all tuples of elements of $A\cup\{w\}$, and inductively for lower
indices. The case of $S^{0}$ is an edge case, but the definition
of good sequence still ensures that $S^{0}(x,y)\equiv[(x,y)^{*}\in S^{0}A]$.
\item For every other triple of elements from $V_{\alpha+1}$, set $S_{\alpha+1}^{n+1}(\vec{x})$
to hold true.
\end{itemize}
Note that the first and second item in the definition do not contradict
each other: If a tuple $\vec{x}$ fits both bullet points, then either
all its elements are in $A$, in which case the definition of good
sequence ensures that there is agreement between both possible definitions
of $S_{\alpha+1}^{n+1}(\vec{x})$, or one of its elements is $w=w^{n+1}(\beta,mA,A,\dots,S^{n+1}A)$,
in which case $S_{\gamma}^{n+1}(\vec{x})$ would have been defined
to agree with the second bullet point for $\gamma$ the smallest ordinal
such that $w\in V_{\gamma}$.

This induces a well-defined relation $S^{n+1}$, whose main defining
properties are as follows:
\begin{lemma}
\label{lem:sneg}Provably in $\ZFniptc$, $\neg S^{n}(\vec{x})$ iff
$(\forall y)S^{n+1}(\vec{x},y)$.
\end{lemma}
\begin{proof}
By construction.
\end{proof}
%
\begin{lemma}
\label{lem:sflex_general}Let $M$ be a model of $\ZFniptc$. Let
$A$ be a finite subset of $M$, and suppose that we wish to find
an element $w$ that relates to all elements of $A$ with regard to
the predicates $S^{0},\dots,S^{n+1}$ in a prescribed manner. The
following is a sufficient condition for there to exist an element
$w$ of $M$ satisfying this prescription: For every $0\leq i\leq n$,
the prescription satisfies the rules
\begin{itemize}
\item The prescription agrees with the pre-existing relations on tuples
that don't contain any instance of $w$,
\item For all $x\in A\cup\{w\}$, it is prescribed that $S^{0}(x,w)$ and
$S^{0}(w,x)$.
\item For all $(\vec{x},y)\in(A\cup\{w\})^{i+2}$ such that $\neg S^{i}(\vec{x})$
is prescribed, it is also prescribed that $S^{i+1}(\vec{x},y)$.
\end{itemize}
\end{lemma}
\begin{proof}
By construction.
\end{proof}
These two properties are enough to establish:
\begin{thm}
\label{thm:gen_pakhomov_jump_inversion}Let $D$ be an $X'$-computable
model of $T^{n}$ in the predicates $S^{i},\dots,S^{n}$, with $i<n$.
Then, the reduct $D\upto(S^{i+1},\dots,S^{n})$ admits an $X$-computable
copy.
\end{thm}
\begin{proof}
We will apply \thmref{main_final} to $D$ which, as an $X'$-computable
model over a finite signature, is also an $X'$-computable $X$-c.e.-typed
model. We need only establish that $D$ has the QETP over $\Lang=(S^{i+1},\dots,S^{n})$
(plus negations) and $\Lang'=(S^{i},\dots,S^{n})$ (plus negations).
\global\long\def\abs#1{\left|#1\right|}%

We construct $\chi$, $\tau$, and $\etau$:
\begin{itemize}
\item First, we construct $\chi_{q}$ for $q(\vec{x},\vec{y})$ a quantifier-free
$\Lang$-type with all variables distinct.

We start by considering the case where $q(\vec{x},y)$ only has one
$y$-variable. In this case, the formula $(\exists y)q(\vec{x},y)$
is true if, and only if, there exists an element $y$ satisfying a
certain prescription of predicates $S^{i+1},\dots,S^{n}$. We show
that the following three statements are equivalent:
\begin{itemize}
\item $P(\vec{x})\equiv(\exists y)q(\vec{x},y)$
\item $Q(\vec{x})$: For every instance of $\neg S^{j+1}(\vec{z}_{0},z_{1})$
in $q(\vec{x},y)$:
\begin{itemize}
\item If $j>i$, we require that $S^{j}(\vec{z}_{0})$ be in $q$, and\footnote{Note: This demand is happening outside of first-order logic. Rather,
it's a check that we metatheoretically make of $q(\vec{x},y)$, and
if this check is failed, we set $Q(\vec{x})\equiv\bot$.}
\item If $j=i$, $z_{1}\equiv y$, and $\vec{z}_{0}$ is a subtuple of $\vec{x}$,
we require $S^{i}(\vec{z}_{0})$,
\end{itemize}
Moreover, we also demand $(q\upto\vec{x})(\vec{x})$. In other words,
whatever demand $q$ makes of the tuple $\vec{x}$, we also demand
it.
\item $R(\vec{x})\equiv$ There is a $y$ such that $q(\vec{x},y)$ and
the relations $S^{0}$ to $S^{i}$ hold for all tuples including elements
from $(\vec{x},y)$ that include at least one instance of $y$.
\end{itemize}
Indeed, $P(\vec{x})\rightarrow Q(\vec{x})$ follows from the fact
that $(\forall\vec{z}_{0})((\exists z_{1})\neg S^{j+1}(\vec{z}_{0},z_{1}))\rightarrow S^{j}(\vec{z}_{0})$
(and the assumption that $q$ is a type), $Q(\vec{x})\rightarrow R(\vec{x})$
follows by \lemref{sflex_general}, and $R(\vec{x})$ obviously implies
$P(\vec{x})$.

Now, still in the case where there is only one $y$-variable, we set
$\chi_{q}(\vec{x})\equiv Q(\vec{x})$ as above. We will now handle
the case where there are two $y$-variables, say $y_{1}$ and $y_{2}$,
and an obvious induction will handle the general case. We will use
the formulas $P$, $Q$, and $R$ constructed above for different
formulas $q$. We will use a subscript to indicate which formula $q$
is being used.

Let $q(\vec{x},y_{1},y_{2})$ be a quantifier-free $\Lang$-type.
Then, let us look at $\chi_{q}(\vec{x},y_{1})$. As above, it's either
$\bot$ (if the first bullet point of the definition of $Q_{q}(\vec{x},y_{1})$
fails), in which case the quantifier elimination is trivial, or otherwise
it consists of $q\upto(\vec{x},y_{1})$, together with a few \emph{positive}
requirements of type $S^{i}(\vec{z}_{0})$ with $\vec{z}_{0}$ a subtuple
of $(\vec{x},y_{1})$. Let $q'(\vec{x},y_{1})$ be $q\upto(\vec{x},y_{1})$
together with those requirements that include $y_{1}$, and $q^{*}(\vec{x})$
be the remainder of these requirements. Note that
\[
(\exists y_{1})(\exists y_{2})q(\vec{x},y_{1},y_{2})\leftrightarrow(\exists y_{1})\chi_{q}(\vec{x},y_{1})\leftrightarrow q^{*}(\vec{x})\land(\exists y_{1})q'(\vec{x},y_{1}),
\]
and crucially that, since all $S^{i}$-requirements made by $q'$
are positive, $(\exists y_{1})q'(\vec{x},y_{1})$ is logically between
$P_{q'}(\vec{x})$ and $R_{q'}(\vec{x})$, and as such is equivalent
to both, and as such is equivalent to $Q_{q'}(\vec{x})$. Then, we
conclude $(\exists\vec{y})q(\vec{x},\vec{y})\leftrightarrow q^{*}(\vec{x})\land Q_{q'}(\vec{x})$,
and so we define the latter as $\chi_{q}$ in this scenario. The induction
continues in a similar fashion, in the event that there are more $y$-variables,
with the essential point being that every $S^{i}$-demand being made
at any step is positive.
\item Now, we construct $\tau$. Given $q(\vec{x},y,\vec{z})$, we set $\tau_{q\top}(\vec{x},y)$
to be the demand that every $S^{i}$-relation including $y$ at least
once holds, and that $\chi_{q}(\vec{x})$ holds.
\item Finally, $\etau_{q\top Q}(\vec{x})$ is simply $\chi_{Q}(\vec{x})$.
\end{itemize}
The proof that these formulas are witnesses to the QETP is simple
and follows the same $P$, $Q$, $R$ reasoning used to establish
that $\chi_{q}(\vec{x})\leftrightarrow(\exists y)q(\vec{x},y)$.
\end{proof}
\thmref{nestedpakhomov} immediately follows.

As explained in the introduction, we obtain as a corollary:
\begin{thm}
\label{thm:pakhomov_answer}For every $n$, there is a theory definitionally
equivalent to ``$\mathsf{PA}$ plus all true $\Pi_{n}$ sentences''
that has a computable non-standard model.
\end{thm}

\section{Further Questions}

It is a standard theorem (or definition) that a $\mathrm{PA}$ degree
is the same as one which computes a nonstandard model of $\mathsf{PA}$.
On the other hand, we saw that there is a computable nonstandard model
of the definitionally equivalent theory $T_{0}$ constructed by Pakhomov
-- and so, by a theorem of Knight \cite{knight_degrees_coded_in_jumps},
there exists one in every degree. Let us define, for a theory $T$
that is definitionally equivalent to $\mathsf{PA}$, its \emph{nonstandard
spectrum} as the set of degrees of nonstandard models of $T$.
\begin{question}
If $T$ is a c.e.\ theory that is definitionally equivalent to $\mathsf{PA}$,
we can see that its nonstandard spectrum is at most the set of all
Turing degrees, and at least the set of all $\mathrm{PA}$ degrees.
Does there exist such a theory whose nonstandard spectrum is strictly
in-between these two cases? More generally, what can the nonstandard
spectrum of a c.e.\ theory definitionally equivalent to $\mathsf{PA}$
look like?
\end{question}
The methods used in this paper are too coarse to address this question;
they would have to be cleverly modified to stop the nonstandard spectrum
from including $0$.

We mentioned in the introduction that Lutz and Walsh \cite{lutz_walsh_tennenbaum}
produced a c.e.\ consistent theory such that no definitionally equivalent
theory has a computable model. In fact, they defined, for each infinite
computable binary branching tree $R$ with no ``guessable'' path,
a consistent theory $T_{\mathrm{LW}}(R)$ such that no theory definitionally
equivalent to $T_{\mathrm{LW}}(R)$ has a computable model. We wonder
how close such models are to being computable.
\begin{question}
Define the Lutz--Walsh spectrum of an infinite computable binary
tree $R$ none of whose paths are guessable, say $\mathrm{LWSpec}(R)$,
as the set of Turing degrees that compute a model of a theory definitionally
equivalent to $T_{\mathrm{LW}}(R)$. It is clear that this $\mathrm{LWSpec}(R)$
always contains all $\mathrm{PA}$ degrees, and in \cite{lutz_walsh_tennenbaum}
it is proven that $\mathrm{LWSpec}(R)$ does not contain $0$ (if,
again, none of the paths of $R$ are guessable). Does $\mathrm{LWSpec}(R)$
ever contain any degrees that are not $\mathrm{PA}$? If so, how far
from $\mathrm{PA}$ can we get? Is there such a tree $R$ for which
$\mathrm{LWSpec}(R)$ consists of all non-zero degrees? Is there any
nonzero degree that can never be in $\mathrm{LWSpec}(R)$?
\end{question}
\medskip{}

In the statement of \thmref{nestedpakhomov}, the arity of the predicate
used to replace set inclusion increased by $1$ for every jump in
complexity. We wonder if this is a logical necessity or merely an
artifact of our construction.
\begin{question}
For a given value of $n$, does there exist a theory $T_{Q}^{n}$,
axiomatizing a single ternary predicate $Q(x,y,z)$, such that $T_{Q}^{n}$
is definitionally equivalent to ``$\mathsf{PA}$ plus all $\Pi_{n}$
truths'' which admits a computable nonstandard model? What if we
require the predicate to be binary, instead?
\end{question}
\medskip{}

While \thmref{main_final} guarantees the existence of a $0'$-computable
isomorphism between the original structure and its computable copy,
most results in the literature can only guarantee (out of necessity)
a $0''$ or $0'''$-computable isomorphism. This is explained by the
fact that, in our proofs, we must first create a $0'$-computable
c.e.-typed copy which is isomorphic to the original structure, though
this isomorphism may itself need to be complex. Hirschfeldt asked
whether one could obtain a more direct result, for example by weakening
the assumptions of the QETP.
\begin{question}[Hirschfeldt]
Does there exist a version of \thmref{main_final} that requires
weaker assumptions, possibly using an infinite injury method instead
of finite injury, from which theorems such as \thmref{JIlinear-1}
(jump inversion for linear orders) or \propref{eqvinf} (jump inversion
for equivalence relations with infinitely many infinite classes) may
be obtained more easily?
\end{question}
Many general results in computable structure theory necessitate the
use of a finite relational signature, oftentimes for reasons similar
to those that led us to the notion of $0'$-computable c.e.-typed
structure. We wonder if our definition would be applicable to other
scenarios.
\begin{question}
Are there results in the literature, applicable to finite relational
signatures, which could be further (and usefully) generalized using
the notion of $0'$-computable c.e.-typed structure?
\end{question}
\medskip{}

In Section \ref{subsec:qetp}, we found it useful to work with $\llorc$
formulas, which are related to but not exactly the same as $\Sigma_{1}^{\mathrm{c}}$
formulas. Ash, Knight, Manasse and Slaman showed that the definability
of a relation by a $\Sigma_{1}^{\mathrm{c}}$ formula is related to
computability-theoretic properties, namely the notion of being a r.i.c.e.
(relatively intrinsically c.e.) relation; see \cites[Theorem II.16]{montalban_cst}
for details. We wonder if $\llorc$ formulas admit a similar treatment.
\begin{question}
Is there a computability-theoretic characterization of those relations
in a structure that are definable by a $\llorc$ formula?
\end{question}

\section*{Acknowledgements}

I would like to thank Denis Hirschfeldt, Miles Kretschmer, and Patrick
Lutz for their valuable feedback, help, and advice, without which
this work would not have been possible.

\printbibliography

\end{document}
